\chapter{瓶装快乐}

地铁列车缓缓加速，驶离站台。耳边响起VVVF\footnote{Variable Voltage Variable Frequency，可变电压可变频率。在这里指驱动地铁列车的采用IGBT-VVVF（绝缘栅双极型晶体管-可变电压可变频率）技术的异步或永磁同步牵引电机}的悦耳歌声；可惜不一会儿就听不见了，因为铃和彩名戴上了耳机。这里没有两人共享耳机的桥段，因为现在的手机已经可以同播共享了，即使使用的是各自的耳机。这里也没有因车厢拥挤而两人抱在一起的桥段，因为今天是休息日（普遍意义来说），没有早晚通勤高峰。彩名出奇安静地聆听着音乐，想要随着拍子摇摆双腿，却发现她的腿太长且列车上的座位太低而无法做到，于是只好点着头轻轻拍起膝盖。铃隐约听到她在随着歌曲小声附和着唱词：
\begin{quote}
悲剧结局真没意思\\
还没到让你觉得讨厌的程度啦！
\end{quote}

看到彩名沉浸在音乐中，铃便开始处理自己的事务，也就是检查今早的讯息。她打开手机上的即时通信软件，开始虐杀红点。绝大多数讯息都与她无关；即使有关，也仅仅是单向度的通知，而非请求回复的祈愿。最后只留下一条讯息：
\begin{quote}
{\heiti 太上小君:}~我平时有没有压榨你们\\
{\heiti 太上小君:}~btw彩名在吗也替我问问她\\
{\heiti 太上小君:}~我不知道怎么翻译成她的私人语言yk
\end{quote}

铃搜索了一下“yk”是什么意思，然后开始思考该怎么回复大老板。

\begin{branch}
\goodend 铃立刻回复“怎么会呢，教授您最关心我们了。”手机另一边的大老板看到这样的回复，心花怒放，不久之后帮助铃申请了一个很大的基金。
\end{branch}

看样子大老板是被今天的跳楼事件吓到了。显然，一个真正压榨学生的老板是不会发讯息来问自己有没有在压榨的。而且，如果真要说有压榨，那铃感觉相比之下大老板压榨他自己起来更狠。铃于是回复“没有吧。”她转过头问了问彩名；彩名把耳机取下来后，她又问了一遍，而后回复“彩名也说没有。”

彩名不再听歌后呆呆地坐着，听着VVVF的舒缓的乐音。看来有时候机械的声音比人造的乐声更能浸染人心。铃思考着将VVVF纳入心理治疗的可能性。就这样不知道过了多久，彩名突然开口，轻轻地问出一个问题：“地铁列车驶向哪里呢？”

铃不清楚这是个认识论问题还是本体论问题。她只好回答：“我们在体育场站下车。”

彩名点了点头；看样子她接受了这个回答，尽管这个答案似乎没能回答她的问题。彩名于是提出下一个问题：“地铁列车只能看到隧道，”她望着车窗外的黑暗接着黑暗，“那么列车能够知道天空吗？”

前提是对的，因为这条地铁线路确实没有地上部分。铃仍然没有明白彩名想要问什么，她只能以问题回答问题：“你想说的是休谟的色盲论证吗？”彩名没有回应，铃于是继续说，“康德会说对天空的知识依赖于普遍的范畴而非个体的感官。”

“你知道发现色盲症的道尔顿本身就是色盲吗？”彩名突然扭转话题。

铃点点头：“红绿色盲是X连锁隐性遗传的疾病。”或者说是一种基因疾患，“不过道尔顿没能发现他的疾病与遗传和性别的关系。”

“我们已经得到人类全部的基因组序列了哦！T2T！”铃想起了彩名所说的这个成果，从端粒到端粒的人类基因组测序：“那是个大工程。”

“可喜可贺！人类现在可以被叫做‘人类-3B’了。”

“DNA有四种碱基，”铃指出彩名的错误，“信息量应该乘以二。T2T测序得到人类全基因组有约三十亿碱基对，所以应该有6B的信息。”

彩名承认了自己的失误，接着说：“你不觉得生物体和大语言模型很像吗？只不过前者转录的是密码子，而后者解码的是词元。所以按照LLM的命名方式，我们就叫‘人类-6B’吧！”

铃认为这是个荒谬的比喻；即使承认其相似性，也难以产生实际效用。她按照彩名的方式予以反驳：“人体的每个细胞都包含整套基因组，所以整个人体的参数量远远超过6B。你这样等于是否定了表观遗传学家所有的努力。”铃立刻发现自己的话并不严谨，她补充道，“我是说体细胞；还有，红细胞除外。”

“没错，中心法则现在变成‘中心独断’了呢！”彩名望着车厢内的显示屏；但铃知道她并不在观看其上的视频，而是在她的脑海中遨游，“每个细胞都包含着如何创造人类的全部知识，但它们选择了自我分化，通过不作为，将自己设定为作为它们的世界的人体的一部分。”

“并不是‘选择’，”铃质疑细胞是否有能够作选择的自由意志，“这是由一套复杂的机理决定的，涉及基因选择性表达、表观遗传调控、细胞间信号调控，等等。”

“我没说其中没有一套复杂的机理。就像世界模型的潜空间包含世界上所有的知识，而它能够输出什么仍然依赖复杂的后训练。”彩名突然站起身，“体育场站到了。我们走吧。”她向铃伸出了手。
